\bumppoints{20}

\question
A rock sample contains an isotope that has a half-life of 21 million years. The sample contains 6.25\% of the parent element and 93.75\% of the daughter element. How many millions of years old is the sample?

\begin{multicols}{2}
	\begin{choices}
\choice 21. %dontshuffle 
\choice 42. %dontshuffle 
\choice 63. %dontshuffle 
\correctchoice 84. %dontshuffle 
\choice 105. %dontshuffle 
\choice 126. %dontshuffle 
\end{choices}
	
	\columnbreak
	
	\includegraphics[width=0.45\textwidth]{radiometric_dating}
\end{multicols}

\question
Two species of \textit{Heliconius} butterflies occupy different habitats but the two species occur together in one area where their habitats are most similar. In this area, the two species often interbreed and hybrid offspring have been observed.  Yet, biologists continue to classify the butterflies as two separate species.  Which of the following observations would provide the best support for classifying the butterflies as two separate species? 

\begin{choices}
\choice The hybrids are cryptic species.
\choice The two species are morphological species, not biological species.
\choice The two species are biological species, not morphological species.
\choice The two species are cryptic species.
\correctchoice The hybrids have much lower fitness in the parental habitats.
\choice The two species are reproductively isolated because they occupy very different habitats.
\end{choices}

\question
Which of the following factors would \emph{not} contribute to allopatric speciation?

\begin{choices}
\correctchoice Gene flow between the two populations is extensive.   
\choice Different mutations begin to distinguish the gene pools of the separated populations.   
\choice The separate population is small, and genetic drift occurs.   
\choice One population is exposed to different selection pressures than the other population.   
\end{choices}

\question
A hybrid zone is properly defined as 

\begin{choices}
\choice a zone where the two closely related species remain distinct through hybrid fusion. 
\choice an area where the ranges of two closely related species overlap, but do not interbreed. 
\choice an area where members of two closely related species intermingle, but gene flow is prevented by prezygotic barriers. 
\choice an area where members of two closely related species intermingle, but gene flow is prevented by allopatric barriers. 
\correctchoice an area where mating occurs between members of two closely related species, producing at least some viable offspring. 
\end{choices}

\question
The males of some fishes and birds have bright colors that males use during territorial displays. The existence of these bright colors in unrelated organisms is an example of 

\begin{choices}
\correctchoice convergent evolution. 
\choice homology. 
\choice descent with modification. 
\choice intersexual selection. 
\choice common ancestry. 
\end{choices}

\question
Lakes in east Africa have been suffering pollution, causing a decline in species diversity of cichlid fishes. This decline in diversity is most likely due to

\begin{choices}
\choice hybrid sterility. 
\choice stability. 
\correctchoice fusion. 
\choice hybrid inviability. 
\choice reinforcement. 
\end{choices}

\question
The studies of columbine flowers showed that spurs act to maintain separate \textit{Aquilegia} species by

\begin{choices}
\choice attracting different pollinators. 
\choice offering different rewards to different pollinators. 
\correctchoice influencing pollen removal from and depositing on flowers. 
\choice causing flowers to grow in different habitats. 
\end{choices}

\newpage

\begin{multicols}{2}
	\question
	The figure at right shows the average number of pollen grains remaining on \textit{Aquilegia pubescens} after visits by hawkmoths. Control flowers had normal length spurs. Shortened flowers had spurs that were cut short and tied by the researchers. Based on the figure, tell which, if any, type of reproductive isolation is supported by the data. 
	
	\begin{choices}
\choice Temporal isolation. 
\choice Gametic isolation. 
\choice Habitat Isolation. 
\correctchoice Mechanical isolation. 
\choice The figure does not provide support for the above types of reproductive isolation. %dontshuffle
\end{choices}
	
	\columnbreak
	
	\includegraphics[width=0.45\textwidth]{/Users/goby/pictures/teach/163/case_studies/pollen_grains_remaining}
\end{multicols}

\question
The type of fossil that is formed from tracks left behind by an organism is called

\begin{choices}
\choice chemical. 
\choice bony. 
\choice permineralized. 
\correctchoice trace. 
\choice pressurized. 
\end{choices}

\question
What does the biological species concept use as the primary criterion for determining whether two populations are the same species?

\begin{choices}
\choice Morphological similarity.   
\choice Ecological differences.   
\choice Geographic isolation.   
\correctchoice Gene flow.   
\choice Allopatric speciation.   
\end{choices}
